\section{Matrix moments and extension across the contact point}
\label{sec:matrix}
\label{mat:section}

The transfer density constructed in Section~\ref{sec:spheres} has a
uniform $L^2$ bound whenever the two normalized principal tensors differ
by at most a small multiple of the sphere radius.  We now improve that
comparison from order $r$ to order $r^{3/2}$.  This strict improvement will
allow the radii to tend to zero while keeping one fixed dilation of the
original coordinates.

The scalar argument uses affine and quadratic moments of the transfer
density; see \cite[Section~6]{ScalarCompanion2026}.  In a harmonic frame,
the first moments are matrices and need not vanish.  We approximate
them by a matrix-valued displacement independent of the shrinking
parameter.  Its component complementary
to the real scalar matrices satisfies a first-order system with
conformal Killing principal part.  The
finite-type property of that system provides the additional power of
the radius.

\subsection{A finite family of harmonic jets}

We use the coordinates and harmonic frames of
Proposition~\ref{bdy:contact}.  Before dilation, their matrix coordinate
functions satisfy
\[
 \frac12\operatorname{Re}\operatorname{tr}X_j^l(x)=x^l,
 \qquad X_j^l(0)=0,
 \qquad \partial_iX_j^l(0)=\delta_i^l I,
 \qquad 1\le l\le n.
\]
There are also finitely many paired harmonic columns whose value and
gradient vanish at the origin and whose first-side Hessians form a
basis of the vector-valued trace-free symmetric Hessians there.  All
these pairs agree on the known side of the contact hypersurface.

For the dilation parameter $\lambda$, replace the coordinate matrices
by
\begin{equation}
 \widehat X_j^l(y)=\lambda^{-1}X_j^l(\lambda y),
 \label{mat:renormalized-coordinates}
\end{equation}
and replace each of the additional columns $u_j$ by
$\lambda^{-2}u_j(\lambda y)$.  Let
\[
 (v_{1,\nu},v_{2,\nu}),\qquad 1\le\nu\le N,
\]
be the resulting list of column pairs, including all columns of the
$\widehat X_j^l$.  The list is finite and fixed independently of the
shrinking parameter $\delta$.  Taylor's formula at the origin shows
that these functions have uniform bounds of every fixed smooth order
on the fixed dilation range.  To obtain the second-side bounds, note
that each paired difference, expressed in the common coordinates,
vanishes on an open set accumulating at the origin.  Every derivative
therefore vanishes at the origin by continuity.  Each additional
column pair thus has zero value and gradient on both sides, as
required for its quadratic rescaling; the coordinate pairs have the
common zero value required for linear rescaling.  Taylor's formula
then gives uniform bounds of every fixed order on both sides.

Recall the harmonic-frame equations, written with positive definite
second-derivative coefficients:
\begin{equation}
 \mathcal L_1=a^{ij}\partial_i\partial_j+B_1^i\partial_i,
 \qquad
 \mathcal L_2=Q^{ij}\partial_i\partial_j+B_2^i\partial_i,
 \qquad Q=a+h.
 \label{mat:operators}
\end{equation}
The tensors $a,Q$ are real and scalar in the fiber index; the $B_j^i$
are complex $2\times2$ matrices.  All coefficients have uniform smooth
bounds for small $\lambda$.  There is no zeroth-order term because the
columns of the frames are harmonic.  In ordinary coordinates,
$B_j^i\longrightarrow0$ smoothly under dilation, whereas
$a\longrightarrow e^{2Lt_0}I_n$.

\begin{lemma}[Uniform spanning by the selected jets]
\label{mat:spanning}
For all sufficiently small $\lambda$ and every $y\in Y$, the vectors
\[
 \big(\partial_i v_{1,\nu}(y),
             \partial_i\partial_j v_{1,\nu}(y)\big),
 \qquad 1\le\nu\le N,
\]
span, over $\mathbb C$, the space of pairs $(p,H)$ with
$p_i\in\mathbb C^2$ and
$H_{ij}=H_{ji}\in\mathbb C^2$ satisfying
\begin{equation}
             a^{ij}H_{ij}+B_1^ip_i=0.
 \label{mat:jet-constraint}
\end{equation}
The coefficients realizing prescribed such data may be chosen by a
smooth right inverse, with uniform bounds to every fixed order.
\end{lemma}

\begin{proof}
As $\lambda\to0$, the coordinate columns converge smoothly to
$y^l e_\alpha$, $1\le l\le n$, $1\le\alpha\le2$, where
$e_\alpha$ is the standard basis of $\mathbb C^2$.  The additional
columns converge to the trace-free quadratic polynomials specified by
their Hessians at the origin.  At any point $y$, the latter give
arbitrary trace-free Hessians, and the affine columns correct their
gradients to any prescribed values.  Thus the limiting jets span all
$(p,H)$ with $\sum_iH_{ii}=0$, uniformly for $y\in Y$.

The spaces in \eqref{mat:jet-constraint} form a smooth vector bundle of
constant rank: contraction with the positive definite tensor $a$ is
onto $\mathbb C^2$.  They can be trivialized by the gradient and the
trace-free part of the Hessian, the remaining trace being determined
by \eqref{mat:jet-constraint}.  In this trivialization the selected jets
give a smooth matrix of constant full row rank at $\lambda=0$.
Compactness of $Y$ preserves this rank, with a positive lower bound on
its nonzero singular values, for small $\lambda$.  If this matrix is
$T$, the right inverse $T^*(TT^*)^{-1}$ has the asserted smoothness and
uniform bounds.  Each selected column satisfies
\eqref{mat:jet-constraint}, so its range is precisely the stated space.
\end{proof}

\subsection{Displacement and weighted remainders}

Throughout the next three subsections we fix $0<\delta\le1$ and assume
the provisional comparison
\begin{equation}
 H_8(y):=\sum_{|\alpha|\le8}|\partial^\alpha h(y)|
                    \le\varepsilon r_\delta(y),\qquad y\in Y.
 \label{mat:provisional}
\end{equation}
Write $r=r_\delta$.  Since $r_{\delta'}\ge r_\delta$ when
$\delta'\ge\delta$, the bound \eqref{mat:provisional} also holds with
every intermediate radius used to continue from $\delta'=1$ to
$\delta'=\delta$.  Proposition~\ref{sph:transfer} and
Proposition~\ref{ang:density} give
\begin{equation}
 \begin{split}
 v_{2,\nu}(y)
    &=\int_{S_y}\psi(y,\omega)\,
          v_{1,\nu}\big(\exp_y^\gamma(r(y)\omega)\big)
          \,d\sigma_y(\omega),\\
 \int_{S_y}\psi(y,\omega)\,d\sigma_y(\omega)&=I,
 \qquad \|\chi\psi\|_{L^2(Y\times S_y)}\le C.
 \end{split}
 \label{mat:transfer-input}
\end{equation}
The constant is independent of $\delta$ and of sufficiently small
$\lambda$.  The measure in the last norm is
$dV_\gamma(y)d\sigma_y(\omega)$; on the fixed coordinate range it is
uniformly equivalent to the corresponding Euclidean base measure.

Set $x(y,\omega)=\exp_y^\gamma(r(y)\omega)$.  The first Taylor
coefficients of the transfer formula involve the matrix moments
\[
 M^i(y)=\int_{S_y}\psi(y,\omega)
                   \big(x^i(y,\omega)-y^i\big)\,d\sigma_y(\omega).
\]
These moments depend on $\delta$ through the density and the sphere.
We need a displacement that remains fixed as the radii shrink.  The
paired coordinate matrices provide such a choice, and the moment
estimate below will show that it differs from $M$ only by a
second-order error.

Define matrices $D^i(y)$, $1\le i\le n$, by the linear system
\begin{equation}
 \widehat X_2^l-\widehat X_1^l
        =\sum_iD^i\partial_i\widehat X_1^l,
                   \qquad 1\le l\le n.
 \label{mat:displacement}
\end{equation}
The map on the right sends the $n$ matrices $D^i$ to $n$ matrices by
right multiplication with the indicated coefficients.  It converges
uniformly, with derivatives, to the identity map because
$\partial_i\widehat X_1^l\to\delta_i^l I$.  Hence
\eqref{mat:displacement} defines smooth $D^i$ with uniform smooth
bounds.  In particular, $D$ is independent of $\delta$.  For each
selected pair define
\begin{equation}
 R_\nu=v_{2,\nu}-v_{1,\nu}-D^i\partial_i v_{1,\nu}.
 \label{mat:remainder}
\end{equation}
These remainders also have uniform smooth bounds and are independent
of $\delta$.

\begin{lemma}[Weighted moment bounds]
\label{mat:moments}
Under \eqref{mat:provisional},
\begin{equation}
 \|\chi D/r\|_{L^2(Y)}\le C,
 \qquad
 \|\chi R_\nu/r^2\|_{L^2(Y)}\le C,
                   \qquad 1\le\nu\le N.
 \label{mat:weighted-moments}
\end{equation}
The constants are uniform in $\delta$ and small $\lambda$.
\end{lemma}

\begin{proof}
Writing $\rho(y)=\|\psi(y,\cdot)\|_{L^2(S_y)}$, uniform normal
coordinates and Cauchy--Schwarz give $|M(y)|\le Cr(y)\rho(y)$.
Taylor expansion in the fixed chart, followed by
\eqref{mat:transfer-input}, gives
\begin{equation}
 v_{2,\nu}-v_{1,\nu}
      =M^i\partial_i v_{1,\nu}+E_\nu,
             \qquad |E_\nu(y)|\le Cr(y)^2\rho(y).
 \label{mat:taylor-transfer}
\end{equation}
This calculation uses only the uniform second derivatives of the
selected columns and $|x-y|\le Cr$; the coordinate displacement is a
scalar, so no matrix factors have been interchanged.

Apply the same calculation to the coordinate matrices.  Subtracting
\eqref{mat:displacement} shows that the uniformly invertible map
defining $D$ sends $D-M$ to an error bounded by $Cr^2\rho$.  Therefore
\[
 |D-M|\le Cr^2\rho,
 \qquad |D|\le Cr\rho,
 \qquad |R_\nu|\le Cr^2\rho.
\]
Multiplication by $\chi$ and integration prove
\eqref{mat:weighted-moments}.
\end{proof}

The next estimate turns these weighted $L^2$ bounds into bounds for
enough derivatives to compare the equations.  No derivative of the
shrinking radius is taken in this interpolation step.

\begin{lemma}[Pointwise interpolation]
\label{mat:interpolation}
Choose
\[
 d_*=10+\frac n2,\qquad \eta=\frac1{2d_*},\qquad
 p_*>2d_*,\qquad m\ge\lceil10+3d_*\rceil.
\]
There is a fixed $\rho_0>0$ such that, if
$y\in K$ and $r(y)<\rho_0$, then
\begin{equation}
 \sum_{|\alpha|\le10}|\partial^\alpha D(y)|\le Cr(y)^{1/2},
 \qquad
 \sum_{|\alpha|\le10}|\partial^\alpha R_\nu(y)|
                                      \le Cr(y)^{3/2}.
 \label{mat:pointwise}
\end{equation}
The threshold and constants are independent of $\delta$ and small
$\lambda$.
\end{lemma}

\begin{proof}
By Lemma~\ref{sph:profiles}, $r^{1/p_*}$ has a uniform Lipschitz
constant.  Set $s=r(y)^\eta$.  Since
$\eta>1/p_*$, for $r(y)$ sufficiently small every $z\in B(y,s)$
satisfies
\[
 \big|r(z)^{1/p_*}-r(y)^{1/p_*}\big|
       \le C r(y)^\eta\le\tfrac12 r(y)^{1/p_*}.
\]
Consequently $c r(y)\le r(z)\le C r(y)$ on this ball, with constants
depending only on the fixed $p_*$.  Because $\chi=1$ on a neighborhood
of the compact set $K$, the ball lies in that neighborhood after a
further fixed reduction of $\rho_0$.  The weighted bounds yield
\[
 \|D\|_{L^2(B(y,s))}\le Cr(y),\qquad
 \|R_\nu\|_{L^2(B(y,s))}\le Cr(y)^2.
\]

For any smooth vector- or matrix-valued function $U$ with a bounded
$C^m$ norm, Taylor expansion to degree $m-1$ and equivalence of norms
on the finite-dimensional polynomial space give
\begin{equation}
 |\partial^\alpha U(y)|
    \le C s^{-|\alpha|-n/2}\|U\|_{L^2(B(y,s))}
                    +C_m s^{m-|\alpha|},\qquad |\alpha|<m.
 \label{mat:taylor-interpolation}
\end{equation}
Indeed the Taylor remainder has $L^2$ norm at most
$C_m s^{m+n/2}$, and, after rescaling the ball to unit radius, each
coefficient of its Taylor polynomial is bounded by that polynomial's
$L^2$ norm.  This proves the displayed inequality componentwise.

For $|\alpha|\le10$, the first term of
\eqref{mat:taylor-interpolation} applied to $D$ is at most
$Cr(y)^{1-\eta d_*}=Cr(y)^{1/2}$; for $R_\nu$ it is at most
$Cr(y)^{2-\eta d_*}=Cr(y)^{3/2}$.  The second term is at most
$C_m r(y)^{\eta(m-10)}\le C_m r(y)^{3/2}$, taking
$\rho_0\le1$.  This proves \eqref{mat:pointwise}.
\end{proof}

\subsection{The equation for the displacement}

The weighted moments give a preliminary $r^{1/2}$ bound for $D$.
To improve it, we compare the two harmonic equations on the selected
columns.  This comparison separates the scalar tensor $h$ from the
matrix part of $D$.

For a selected first-side column write $w=v_{1,\nu}$ and $R=R_\nu$.
Using $v_{2,\nu}=w+D^k\partial_kw+R$, expand
$\mathcal L_2v_{2,\nu}=0$ and use
$\mathcal L_1w=0$.  In the third derivatives use
\[
 a^{ij}\partial_i\partial_j\partial_kw
   =-(\partial_ka^{ij})\partial_i\partial_jw
     -(\partial_kB_1^i)\partial_iw
     -B_1^i\partial_i\partial_kw.
\]
The resulting identity is
\begin{align}
 &\big[h^{ij}I+a^{li}\partial_lD^j+a^{lj}\partial_lD^i
                           +\mathcal B^{ij}(D)\big]
                               \partial_i\partial_jw
            +\mathcal Y^i\partial_iw\notag\\
 &\hspace{15mm}=-\mathcal L_2R
       -2h^{ij}(\partial_iD^k)\partial_j\partial_kw
       -D^kh^{ij}\partial_i\partial_j\partial_kw,
 \label{mat:comparison-identity}\\
 &\mathcal B^{ij}(D)
       =-D^k\partial_ka^{ij}
        +\tfrac12\big(B_2^iD^j+B_2^jD^i
                              -D^jB_1^i-D^iB_1^j\big),
 \label{mat:B-definition}\\
 &\mathcal Y^k=B_2^k-B_1^k
       +Q^{ij}\partial_i\partial_jD^k
       +B_2^i\partial_iD^k-D^l\partial_lB_1^k.
 \label{mat:Y-definition}
\end{align}
In particular, $\mathcal Y$ is independent of the selected column.
The order of the factors in \eqref{mat:B-definition} records the
noncommutativity of the matrix drifts.  The drift of $\mathcal L_2$ acting on
$D^k\partial_kw$ places $B_2^i$ on the left of $D^k$, while
differentiating the first-side equation places $B_1^i$ on its right.

Let $\Pi$ be the projection of symmetric contravariant spatial
tensors onto their $\gamma$-trace-free part, acting entrywise on
matrices:
\[
 (\Pi T)^{ij}=T^{ij}
       -\frac1n\big(\gamma_{kl}T^{kl}\big)a^{ij}.
\]
Recall that $\gamma_{ij}h^{ij}=0$.

\begin{lemma}[Projected displacement equation]
\label{mat:projected-equation}
There is a smooth matrix-valued symmetric tensor $\mathcal E$ such
that
\begin{equation}
 h^{ij}I+
 \Pi\big(a^{li}\partial_lD^j+a^{lj}\partial_lD^i
                           +\mathcal B^{ij}(D)\big)
                     =\mathcal E^{ij}.
 \label{mat:projected}
\end{equation}
For $y\in K$ with $r(y)<\rho_0$,
\begin{equation}
       \sum_{|\alpha|\le8}|\partial^\alpha\mathcal E(y)|
                                      \le Cr(y)^{3/2}.
 \label{mat:E-bound}
\end{equation}
\end{lemma}

\begin{proof}
Denote the right-hand side of
\eqref{mat:comparison-identity} for the column $v_{1,\nu}$ by
$Z_\nu$.  Its derivatives through order eight are bounded by
$Cr^{3/2}$ on the stated set.  For $\mathcal L_2R_\nu$ this uses
derivatives of $R_\nu$ through order ten in
\eqref{mat:pointwise}.  Every derivative of either remaining product
contains a derivative of $h$ of order at most eight and a derivative
of $D$ of order at most nine or eight, respectively.  Thus
\eqref{mat:provisional} and \eqref{mat:pointwise} give the factor
$r\,r^{1/2}$.  The other factors are fixed smooth derivatives of the
selected columns.

At each point, Lemma~\ref{mat:spanning} permits linear combinations of
the identities \eqref{mat:comparison-identity} for which the combined
gradient is zero and the combined Hessian is any prescribed
$a$-trace-free symmetric covariant tensor with values in
$\mathbb C^2$.  Such Hessians satisfy
\eqref{mat:jet-constraint} with zero gradient.  They annihilate the
$\mathcal Y$ term and test precisely the $\gamma$-trace-free part of
the coefficient of the Hessian.

More explicitly, choose a smooth uniformly bounded basis $H_s$ of
scalar symmetric Hessians with $a^{ij}(H_s)_{ij}=0$, and use the right
inverse in Lemma~\ref{mat:spanning} to realize $H_se_\alpha$, for
$\alpha=1,2$, with zero gradient.  The resulting combinations of
$Z_\nu$ give the two columns of the contraction of the projected
coefficient with $H_s$.  Inverting this spatial duality defines
$\mathcal E$ and proves \eqref{mat:projected}.  The coefficients
realizing these prescribed jets are smooth functions of $y$ and give
zero combined gradient identically.  We form these combinations
before differentiating, so the $\mathcal Y$ term is identically zero.
The resulting formula expresses $\mathcal E$ using only the $Z_\nu$
and coefficients with uniformly bounded derivatives.  Differentiating
that formula through order eight proves \eqref{mat:E-bound}, without
introducing derivatives of $\mathcal Y$.
\end{proof}

Define the real-linear projection on matrices
\[
 q(Z)=Z-\tfrac12\operatorname{Re}\operatorname{tr}(Z)I,
 \qquad
 \mathcal V_0=\{Z\in\mathbb C^{2\times2}:
                      \operatorname{Re}\operatorname{tr}Z=0\}.
\]
Write $D^i=d^iI+W^i$, where $d^i\in\mathbb R$ and
$W^i\in\mathcal V_0$.  Taking half the real trace in
\eqref{mat:displacement}, and using the exact trace coordinate
identity for both sides, gives
\begin{equation}
 d^l=-\frac12\operatorname{Re}\operatorname{tr}
                    \sum_i W^i\partial_i\widehat X_1^l.
 \label{mat:scalar-direction}
\end{equation}
Thus $d$ is a zero-order real-linear expression in $W$ with uniformly
smooth coefficients.

Apply $q$ in the matrix index to \eqref{mat:projected}.  The term
$h^{ij}I$ vanishes, as do all derivative terms involving $d^iI$.
Substituting \eqref{mat:scalar-direction} in the remaining
zero-order terms gives
\begin{equation}
 \Pi\big(a^{li}\partial_lW^j+a^{lj}\partial_lW^i\big)
             +\mathcal B_0^{ij}(W)=\mathcal E_0^{ij},
 \qquad \mathcal E_0=q(\mathcal E).
 \label{mat:CK-system}
\end{equation}
Here $\mathcal B_0$ is a real-linear zeroth-order operator with
uniformly smooth coefficients, and $\mathcal E_0$ satisfies
\eqref{mat:E-bound}.  There are no hidden first derivatives in
$\mathcal B_0$: differentiation of \eqref{mat:scalar-direction}
occurs only in the terms proportional to $I$, which have already
been removed by $q$.

\subsection{Finite type and the improved estimate}

We prove the finite-type statement needed for
\eqref{mat:CK-system}.  In Euclidean space, conformal Killing vector
fields have degree at most two; see, for example,
\cite[Section~3]{Eastwood2005}.  The calculation below establishes the
corresponding injectivity on third derivatives and includes the
uniformity needed for the variable-coefficient equation.

\begin{lemma}[Prolongation of the conformal Killing principal part]
\label{mat:finite-type}
Let $n\ge3$, let $a=\gamma^{-1}$ be a smooth positive definite tensor
on a coordinate domain, and let $V$ be a finite-dimensional real vector
space.  Suppose that a $V$-valued vector field $W$ satisfies
\[
 \Pi\big(a^{li}\partial_lW^j+a^{lj}\partial_lW^i\big)
                          +C^{ij}(W)=E^{ij},
\]
where $C$ is a smooth real-linear zeroth-order operator.  Every third
derivative of $W$ is then a smooth linear combination of derivatives
of $W$ through order two and derivatives of $E$ through order two.
The coefficient bounds are uniform when $a$ is uniformly elliptic
and the indicated background coefficients range in bounded smooth
sets.
\end{lemma}

\begin{proof}
Differentiate the equation twice.  At a point, its highest-order part
is a linear map from the tensor of third derivatives of $W$ to the
twice-differentiated trace-free symmetric gradient.  We first show
that this map is injective.  A constant linear change of coordinates
sets $a^{ij}=\delta^{ij}$ at the point, and each real component of
$V$ can be considered separately.

A tensor in the kernel defines a homogeneous cubic vector field $Z$
whose trace-free symmetric gradient vanishes: that gradient is
homogeneous quadratic and all of its second derivatives are zero.
Hence
\begin{equation}
 \partial_iZ_j+\partial_jZ_i=2\delta_{ij}\sigma,
                 \qquad \sigma=\frac1n\operatorname{div}Z.
 \label{mat:CK-cubic}
\end{equation}
Differentiating and combining the three permutations of the indices
gives
\[
 \partial_i\partial_jZ_k
   =\delta_{jk}\partial_i\sigma
      +\delta_{ik}\partial_j\sigma
      -\delta_{ij}\partial_k\sigma.
\]
Taking divergence in $k$ yields
\[
 (n-2)\partial_i\partial_j\sigma=-\delta_{ij}\Delta\sigma.
\]
Its trace gives $(2n-2)\Delta\sigma=0$; since $n\ge3$, the Hessian of
$\sigma$ is zero.  The preceding formula therefore gives
$\partial_i\partial_j\partial_lZ_k=0$, so the homogeneous cubic $Z$
vanishes.  This proves injectivity.

Let $T_a$ denote this injective map on third-derivative tensors, using
fixed Euclidean inner products on the finite-dimensional tensor
spaces.  It has the smooth left inverse
$(T_a^*T_a)^{-1}T_a^*$.  Uniform ellipticity and compactness of bounded
coefficient ranges give uniform bounds for this inverse.  All other
terms in the twice-differentiated equation contain at most two
derivatives of $W$ or of $E$.  Applying the left inverse proves the
claim, including the smooth coefficient bounds after further fixed
numbers of differentiations.
\end{proof}

\begin{proposition}[Strict improvement of the principal tensors]
\label{mat:improvement}
With the fixed choices in Lemma~\ref{mat:interpolation}, there are
constants $C$ and $\rho_0>0$, uniform for sufficiently small
$\lambda$, such that \eqref{mat:provisional} implies
\begin{equation}
                 H_8(y)\le C r_\delta(y)^{3/2}
       \qquad\text{if }y\in K\text{ and }r_\delta(y)<\rho_0.
 \label{mat:improved-h}
\end{equation}
On the same set, the derivatives of $D$ through order nine are bounded
by $Cr_\delta^{3/2}$.
\end{proposition}

\begin{proof}
Apply Lemma~\ref{mat:finite-type} to
\eqref{mat:CK-system} with $V=\mathcal V_0$.  For each multi-index
$\alpha$ of order three, it gives an identity of the form
\begin{equation}
 \partial^\alpha W
       =\sum_{|\zeta|\le2}C_{\alpha\zeta}\partial^\zeta W
         +\sum_{|\zeta|\le2}F_{\alpha\zeta}
                                       \partial^\zeta\mathcal E_0,
 \label{mat:prolongation}
\end{equation}
with uniformly smooth coefficient maps.

Fix $y\in K$ with $r(y)<\rho_0$.  By the core geometry of
Section~\ref{sph:core-cutoff}, the vertical segment
\[
 z(s)=(y',s),\qquad -\tfrac12\le s\le y_n,
\]
stays in $K$, and $r(z(s))\le r(y)$.  At its initial point, $D$ and
all its derivatives vanish: the paired coordinate matrices agree on
an open neighborhood in the known lower region.  The same holds for
$W$.

Let $\mathcal W(s)$ consist of all coordinate derivatives of $W$
through order two at $z(s)$.  Differentiating this vector with respect
to $s$ either gives another one of its components or a third
derivative, to which \eqref{mat:prolongation} applies.  Thus
\[
 \mathcal W'(s)=A(s)\mathcal W(s)+F(s),
 \qquad \mathcal W(-\tfrac12)=0,
 \qquad |A(s)|\le C,
 \qquad |F(s)|\le Cr(y)^{3/2}.
\]
The bound for $F$ follows from \eqref{mat:E-bound}, since the entire
segment has radius below $\rho_0$.  Its length is uniformly bounded,
so the elementary integral form of Gronwall's inequality gives
\[
       \sum_{|\alpha|\le2}|\partial^\alpha W(y)|
                                             \le Cr(y)^{3/2}.
\]

This argument holds at every point of the stated set.  Differentiating
\eqref{mat:prolongation} $k-3$ times, for $3\le k\le9$, expresses
the derivatives of $W$ of order $k$ in terms of its derivatives
through order $k-1$ and derivatives of $\mathcal E_0$ through order
$k-1$.  Induction and \eqref{mat:E-bound} therefore give the same
bound through order nine.  Notice that the last step uses precisely
eight derivatives of $\mathcal E_0$.

Equation~\eqref{mat:scalar-direction} now gives the same estimates for
$d$ and hence for $D$.  Finally, differentiate
\eqref{mat:projected} through order eight.  Its right-hand side uses
derivatives of $D$ through order nine and of $\mathcal E$ through
order eight, all bounded by $Cr^{3/2}$.  This proves
\eqref{mat:improved-h}.
\end{proof}

\subsection{One dilation for all shrinking radii}

We have proved the strict improvement assuming the provisional bound.
We now remove that assumption for all positive shrinking parameters.
The order of choices matters: the threshold for the improvement must
be fixed before the final choice of the dilation.

\begin{proposition}[Local equality of the normalized equations]
\label{mat:local}
In the contact coordinates and harmonic frames of
Proposition~\ref{bdy:contact}, there is a neighborhood of the origin
on which the two metrics are smoothly conformal and the
harmonic-frame equations, multiplied by the
positive scalars making their principal tensors equal, coincide.
All the finitely many selected pairs of harmonic-frame columns agree
on that neighborhood.
\end{proposition}

\begin{proof}
First fix the reference geometry and uniform smooth bounds for all
sufficiently small dilations.  Choose $p_*$ large enough for the sphere
geometry, the angular estimate, and
Lemma~\ref{mat:interpolation}.  Next choose $\varepsilon$ sufficiently
small and $R_*$ sufficiently small to satisfy those results, including
$r\beta^2\le\varepsilon$.  As explained in
Remark~\ref{sph:choices}, every fixed-separation range needed for the
continuation and cutoff bounds is then fixed before $\lambda$.
Proposition~\ref{ang:density} and all estimates of this section have
constants uniform for sufficiently small $\lambda$.

Choose $0<\rho_*<\rho_0$, with $\rho_*\le1$, so small that
$C\rho_*^{1/2}\le\varepsilon/2$ in
\eqref{mat:improved-h}.  Thus, wherever that estimate applies and
$r_\delta<\rho_*$, it improves
\eqref{mat:provisional} to
$H_8\le\varepsilon r_\delta/2$.

As $\lambda\to0$, $h\to0$ in $C^8(Y)$.  The numbers
\[
 \rho_*,\qquad \min_Y r_1,\qquad R_*(2t_b)^{p_*}
\]
are positive and fixed independently of $\lambda$.  We can therefore
choose one sufficiently small $\lambda$ such that
\begin{equation}
 \sup_Y H_8\le\frac{\varepsilon}{2}
       \min\big\{\rho_*,\min_Yr_1,R_*(2t_b)^{p_*}\big\}.
 \label{mat:final-dilation}
\end{equation}
We also make the known lower region contain $\{y_n<-1/4\}$ on the
fixed ranges, as in the sphere construction.  Fix this $\lambda$ for
the rest of the proof.

Let
\[
 \mathcal I=\{\delta\in(0,1]:
              H_8(y)\le\varepsilon r_\delta(y)
                         \text{ for every }y\in Y\}.
\]
It contains $1$ by \eqref{mat:final-dilation}.  For any
$\delta\in\mathcal I$, the preceding estimates improve the inequality
to $H_8\le\varepsilon r_\delta/2$ throughout $Y$.  Indeed,
Proposition~\ref{mat:improvement} handles the points in $K$ with
$r_\delta<\rho_*$.  Equation~\eqref{mat:final-dilation} handles all
points with $r_\delta\ge\rho_*$.  At a point outside $K$, either
$y_n<-1/2$, where equality of the coefficients is already known, or
$t>2t_b$.  In the latter case
$\ell_\delta(t)\ge t$, so
$r_\delta\ge R_*(2t_b)^{p_*}$ and
\eqref{mat:final-dilation} again applies.

The set $\mathcal I$ is relatively closed in $(0,1]$ by continuity
and compactness of $Y$.  It is also relatively open: at any positive
$\delta$, the radius has a positive minimum on $Y$, so the strict
factor $1/2$ just proved persists for nearby parameter values in the
weaker defining inequality.  Since $(0,1]$ is connected,
$\mathcal I=(0,1]$.

On the fixed open set $K^\circ\cap\{t<0\}$, which contains the
origin, $r_\delta\to0$ as $\delta\to0$.  The coefficients $h,D,R_\nu$
are independent of $\delta$.  The provisional bound and
\eqref{mat:pointwise} therefore imply
\[
                   h=0,\qquad D=0,\qquad R_\nu=0
\]
there.  In particular all selected pairs agree.  Subtracting their
equations now gives
\[
                 (B_2^i-B_1^i)\partial_i v_{1,\nu}=0
                         \qquad(1\le\nu\le N).
\]
The gradient projection of the space
\eqref{mat:jet-constraint} is all of $(\mathbb C^2)^n$: for any
gradient one can choose the trace of the Hessian to satisfy that
constraint.  Lemma~\ref{mat:spanning} therefore implies $B_1=B_2$.
Together with $Q=a$, this proves equality of the normalized
harmonic-frame operators.  The definition of $Q$ shows that the two
inverse metrics are positive scalar multiples, so the metrics are
smoothly conformal.  Undoing the fixed dilation proves the stated
conclusion in the original contact coordinates.
\end{proof}

The remaining conformal factor and the metric carried by the harmonic
frames are addressed next.  The local conclusion just proved supplies
the equality of equations needed to recover both the Riemannian metric
and a unitary connection identification.
